% 第 25 章　光原来是电磁波
\chapter{光原来是电磁波}

\step{反常识开场}

\begin{mainline}
闭上眼，面朝太阳站一会儿，眼皮都能感到暖意。这很平常，平常到我们从不追问：你和太阳之间隔着一亿五千万千米的空——是真正的空，没有空气，没有水，没有绳子，连尘埃都几乎没有。可是有什么东西从那边出发，穿过这片什么都没有的空间，跑了八分多钟，抵达你的皮肤，把它晒热。

再想想这些事：收音机在房间里收到几百千米外的电台；Wi-Fi 信号穿过墙壁；医院的 X 光穿过你的身体在底片上留下骨头的影子；遥控器一按，电视就换了台；消防员的热成像仪在浓烟里“看见”被困者身上的热。这些“信号”穿过真空、穿过墙壁、穿过身体。它们究竟是什么？在“什么都没有”的地方，它们靠什么前进？

还有一个更隐蔽的悬念藏在客厅里。按下开关，灯“立刻”亮了；打电话，声音“立刻”传到千里之外。这两个“立刻”真的是立刻吗？如果是，那是什么东西跑得无限快；如果不是，它到底有多快，又是什么在跑？

这一章会给出一个出乎意料、却被无数实验反复确认的答案：光、无线电波、X 射线，全都是同一个东西——一份在空间里自己维持自己、向前飞奔的电场与磁场的联合扰动。太阳晒热你的，和收音机收到的，是同一家族的成员。
\end{mainline}

\step{先猜一猜}

\begin{mainline}
在往下读之前，先写下你的三个猜测。

第一，光从太阳到地球，要不要花时间？如果要，要走多久？日常经验里，开灯的瞬间房间就亮了，好像光是“立刻到达”的。

第二，我们已经知道，水波是水在振动，声波是空气在振动（第 17 章）。如果光也是一种波，那么穿过真空的光，是什么东西在振动？

第三，电、磁、光——你从摩擦起电、指南针、阳光那里分别认识的三位老相识——它们之间有没有关系？

关于第一个问题，历史上几乎所有人都猜“立刻到达”。伽利略曾经设计过一个诚实的实验：两个人各提一盏遮住的灯，站在相隔一两千米的山顶上，甲掀开灯罩，乙看见光就掀开自己的灯罩，甲记下从掀罩到看见回应的时间。结果当然什么也测不出来——不是光不花时间，而是光跑完这一两千米只需百万分之几秒，人的反应速度根本跟不上。

真正的突破口来自天上。1676 年，丹麦天文学家罗默注意到一个怪事：木星的卫星钻进木星阴影的时刻，当地球在轨道上远离木星时，总比预计的“迟到”一些，累计可差二十来分钟。他大胆地解释：不是卫星不守时，而是光要走完那段多出来的路程，需要时间。光速不是无限大，只是大得离谱，日常距离上感觉不出来。这个猜测后来被越来越精确的测量证实：光跑一秒，能绕地球赤道七圈半。

第二个和第三个猜测先留在心里。不过可以先把可能的答案摆出来：关于第二个问题，你可能会猜“也许真空里其实有某种看不见的东西在振”，也可能猜“也许有的波根本不需要东西在振”——后一个猜测听起来自相矛盾，却是本章的主角。关于第三个问题，历史上有过一个诱人的线索：光速的数值，居然能从两个纯电磁学实验的常数里算出来。如果电、磁、光真的互不相干，这种巧合没有理由存在。下面我们亲手制造一点“看不见的光”。
\end{mainline}

\step{动手实验}

\begin{experiment}[用一节电池“发电报”]
材料：一台能收中波（AM）的收音机、一节 5 号电池、一小段粗导线。

把收音机调到两个电台之间的空频率，让喇叭只发出沙沙声。把导线一端按在电池负极上，另一端在电池正极上快速地刮擦、通断。你会立刻听到收音机里响起“咔哒、咔哒”的声音——哪怕收音机在几米之外，中间没有任何连线。

发生了什么？导线通断的一瞬间，电流猛然变化，火花附近的电场和磁场跟着猛然变化，一小团电磁波就被“甩”了出去，飞过整个房间，推动收音机天线里的电子轻轻动了一下，电路把这微小的推动放大成“咔哒”声。

你刚刚重复的是科学史上最著名的实验之一：1887 年，赫兹用更大的火花和更认真的测量，第一次亲手造出并接收了电磁波，证实了麦克斯韦二十多年前的纸上预言。你桌上的这套装置，就是他实验的家庭缩小版。

追加一步：把收音机用铝箔整个包起来，再刮擦导线——“咔哒”声消失了。铝箔并没有“吸走”什么，原因在本章后面揭晓。

这个实验里藏着本章全部的主角：变化的电流（源）、飞过房间的“某种东西”（电磁波）、被推动的电子（接收）。接下来我们要回答的是中间那一环——那“某种东西”不借助空气、不借助导线，究竟是靠什么法则穿过房间的。
\end{experiment}

\begin{experiment}[两片偏振片里的黑暗]
找两副偏振太阳镜（或者从报废的液晶屏幕上拆两片偏振膜）。把两片叠起来对着灯，慢慢旋转其中一片：透过的光越来越暗，转到九十度附近几乎全黑。

现在做一件怪事：把第三片偏振片斜着四十五度插进那两片已经“全黑”的组合中间——黑暗处竟然又亮了！

如果偏振片只是“挡光板”，多加一片绝不可能让光变多。先记下这个现象，本章后面会用它撬开一个关键结论：光的振动方向，垂直于它的前进方向。
\end{experiment}

\step{旧模型遇到麻烦}

\begin{mainline}
回到十九世纪中期，物理学家手里的“官方图景”是这样的。

光是波，这一点证据很硬：光会干涉、会衍射（第 35 章细讲），而只有波才会干这些事。牛顿本人曾倾向于把光看作一束微粒，这能干净地解释直线传播和反射；但杨氏让两束光叠出明暗相间的条纹之后，微粒说就撑不住了——两股粒子流叠在一起只会更多，不会“相加反而变暗”。波动说赢了，可它立刻背上一个大包袱：波总要回答“谁在振动”。水波是水在动，声波是空气在动，光穿过真空，谁在动？物理学家只好发明一种叫“以太”的东西：它充满整个宇宙，又必须极其稀薄——行星在其中穿行亿万年毫无阻力；又必须极其坚硬——才能托起光那么高频率的振动。一种物质同时“稀薄到完全感觉不到”又“坚硬到超乎想象”，这两个要求怎么凑都不像真话，越研究越别扭。

另一边，电和磁也各自为战。电荷隔空相吸相斥，磁针隔空偏转，主流的说法是“超距作用”：力瞬间到达，中间不需要任何过程。这个说法用在静电上似乎也够用——库仑定律算出的力分毫不差，谁还管它是“怎么到达”的？可是电磁感应（第 23 章）已经在敲警钟：磁铁插进线圈，线圈立刻有电流，仿佛有什么东西抢先一步等在线圈里。法拉第不接受超距作用，他想象空间里到处是看不见的“场线”，力是场线一根一根传过去的，变化也要花时间沿着场线走。多数同行觉得这只是不懂数学的人的图画。麦克斯韦看出了图画里的真东西，把它翻译成精确的数学语言，并补齐了关键的一块补丁——位移电流：变化的电场本身就是电流的一位“替身”，同样能卷起磁场（第 24 章）。

结果，方程自己开口说话了。它说出一句没人料到的话：在完全没有电荷、没有电流的真空里，仍然存在一种“会自己走路的场”。
\end{mainline}

\step{发明新概念}

\begin{mainline}
麦克斯韦问：真空区域里，他的四条方程允许什么样的电场和磁场？

先回忆一下四条方程各自的意思（第 24 章）：电荷是电场的源头；没有找到过孤立的磁荷；变化的磁场会卷起旋涡状的电场；电流——以及麦克斯韦补上的“变化的电场”——会卷起旋涡状的磁场。在导线和电荷附近，这些场被源拴着；可是在远离一切源的真空里，第三条和第四条互相咬住了对方：能让电场随时间变化的，是磁场的旋涡；能让磁场随时间变化的，是电场的变化。于是麦克斯韦继续追问：这两条能不能自己维持下去，不靠任何源？

答案惊人，而且是肯定的。设某处的电场沿竖直方向，大小沿水平方向有变化；那么第三条和第四条合在一起强制要求：这样的电场必须随时间变化，而它的变化必然伴随一个磁场，这个磁场既垂直于电场、也垂直于传播方向；磁场的分布又随时间整体向前推移。解出来的完整图案是：电场和磁场像两堵互相垂直的波墙，手挽着手，以固定的速度向远方平移——不需要电线，不需要空气，也不需要以太，真空就是它的全部舞台。

这里要避开一个流行的错误说法：“电场先产生磁场，磁场再产生电场，循环接力向前跑。”不是这样。电场和磁场并非一先一后把对方“生”出来，它们在同一份解里同时存在、步调一致；是这份联合图案的空间分布随时间整体向前移动，整个过程受麦克斯韦方程这一整条规律的约束。它像体育场里的人浪：每个人只在原地上起下坐，“浪”却绕场飞奔。它像人浪的地方在于：传播的是状态的变化，而不是物质的迁移；它不像人浪的地方在于：人浪需要看台上一排排的人，电磁波背后没有任何“人”——是场本身在动。

\begin{center}
\begin{tikzpicture}[>=Stealth,scale=0.85]
  \draw[->] (-0.2,0) -- (10.6,0) node[right]{传播方向};
  \draw[->] (0,-1.7) -- (0,2.0) node[above]{电场 $E$};
  \draw[blue!70!black,thick,domain=0:10,samples=120] plot (\x,{1.25*sin(108*\x)});
  \foreach \x in {0.4,1.3,2.2,3.1,4.0,4.9,5.8,6.7,7.6,8.5,9.4}
    \draw[->,blue!50] (\x,0) -- (\x,{1.25*sin(108*\x)});
  \foreach \x in {0.83,4.17,7.5}
    \node[red!70!black] at (\x,-0.42) {\small$\odot$};
  \foreach \x in {2.5,5.83,9.17}
    \node[red!70!black] at (\x,-0.42) {\small$\otimes$};
  \node[red!70!black,anchor=west] at (7.6,-1.35) {\small 磁场 $B$（$\odot$ 出纸面，$\otimes$ 入纸面）};
\end{tikzpicture}
\end{center}

更惊人的还在后面。麦克斯韦从方程算出这份扰动的传播速度，发现它只由两个常数决定：一个来自静电实验，一个来自静磁实验。代入当时的测量值，得到约每秒三十一万千米。而当时最好的光速测量——斐索在 1849 年用旋转齿轮测出的——是每秒三十一万五千千米上下。两个来自完全不同领域的数字几乎重合。麦克斯韦在 1865 年写道：如此吻合，我们有充分理由断定，光本身就是一种电磁扰动。

光、电、磁，三位老相识原来是一家人。

预言写在纸上，还差一只实验的手。1887 年，赫兹在实验室里让感应线圈在两个铜球之间打出火花——火花意味着电流在极短时间内剧烈变化。按照麦克斯韦的方程，这应该向外甩出电磁波。几米外，他放一个开着一道细缝的金属圆环：果然，圆环的缝隙里跳出了微小却清晰的火花。没有导线，没有接触，能量自己飞过了实验室。赫兹还让这些波反射、干涉、偏振，一项一项对上波的全部特征，并测出它的速度与光速一致。麦克斯韦去世八年后，他的方程被实验证明是对的。十几年后，马可尼把同样的火花变成了横跨大西洋的无线电报——从一条方程到一个产业，只隔了一代人的时间。
\end{mainline}

\step{一句话模型}

\begin{oneline}
光就是电场与磁场联合组成的一份扰动：它不需要任何介质，以每秒约三十万千米的速度在真空中自我维持地向前传播；无线电波、微波、红外线、可见光、紫外线、X 射线、伽马射线全是它的家人，区别只在每秒振动多少次。
\end{oneline}

\step{数学翻译器}

\begin{mainline}
先看这个速度是怎样从方程里“长”出来的。两个常数登场：真空介电常数 $\varepsilon_0\approx 8.85\times 10^{-12}$，来自库仑定律，可以理解为在真空中“撑起”一份电场有多费劲；真空磁导率 $\mu_0 = 4\pi\times 10^{-7}\approx 1.26\times 10^{-6}$，来自电流的磁效应，描述电流在真空中“撑起”磁场的本领。麦克斯韦方程给出的波速是
\[
  c=\frac{1}{\sqrt{\varepsilon_0\mu_0}}\,.
\]
代入数字：乘积 $\varepsilon_0\mu_0\approx 1.11\times 10^{-17}$，开平方得约 $3.34\times 10^{-9}$，取倒数得
\[
  c\approx 3.0\times 10^{8}\ \text{m/s}\,.
\]
一组完全从静电、静磁实验里测出的常数，组合起来竟然等于光速。这是物理学史上最漂亮的“撞车”之一：两个互不相干的实验领域，在同一个数字上相遇。

立刻检查这个公式。如果 $\mu_0$ 趋于零——意思是变化的电场不再伴随任何磁场——分母为零，$c$ 变成无穷大，波动瓦解成瞬时的超距作用；$\varepsilon_0$ 趋于零同理。可见两个常数缺一不可，正是“电”和“磁”两只手缺一不可。反过来，$c$ 也绝不可能是零，因为两个常数都是正的实数。方向反转检查：把电场整体反向，磁场也跟着整体反向，传播方向不变——这正是一份持续向前输送的能量该有的样子。

接下来是描述一切波动都少不了的关系：
\[
  c=\lambda f\,,
\]
波长乘频率等于波速。用极端情况检查它。频率 $f$ 趋于零时，波长 $\lambda$ 趋于无穷大——几乎不随时间变化的场“传不动”，退回成第 19、22 章里那种静态场，这正合理。反过来频率极高时波长极短，短到可以钻进原子的缝隙，这正是 X 射线能拍出骨头、伽马射线能穿透钢板的原因。再用生活里的数字检查一遍：调频广播标着 100 兆赫兹，算得波长约 3 米，和你家收音机的拉杆天线同一尺度； Wi-Fi 用 $2.4\times10^{9}$ 赫兹，波长约 12 厘米，正好和微波炉里那团波的波长一样——它们本来就是邻居频段。再看可见光：取波长 $\lambda\approx 500$ 纳米，算出频率 $f = 3\times10^{8}/(5\times10^{-7}) = 6\times10^{14}$ 赫兹，每秒振动约六百万亿次。难怪谁也没法直接“看见”电场在振动——我们只能探测它对物质的平均效果，比如发热、感光、推动天线里的电子。

把频率从低到高排开，就得到电磁波的全家福。同一条方程、同一个速度，只是每秒振动的次数不同：
\end{mainline}

\begin{center}
\begin{tikzpicture}[>=Stealth]
  \fill[blue!10]   (0,0)     rectangle (4.5,0.55);
  \fill[cyan!15]   (4.5,0)   rectangle (6.75,0.55);
  \fill[green!15]  (6.75,0)  rectangle (8.7,0.55);
  \fill[yellow!50] (8.7,0)   rectangle (8.925,0.55);
  \fill[orange!20] (8.925,0) rectangle (10.5,0.55);
  \fill[red!15]    (10.5,0)  rectangle (12,0.55);
  \fill[violet!18] (12,0)    rectangle (13.5,0.55);
  \draw (0,0) rectangle (13.5,0.55);
  \foreach \x in {4.5,6.75,8.7,8.925,10.5,12} \draw (\x,0) -- (\x,0.55);
  \node at (2.25,0.85) {\small 无线电波};
  \node at (5.6,0.85) {\small 微波};
  \node at (7.7,0.85) {\small 红外线};
  \node at (9.7,0.85) {\small 紫外线};
  \node at (11.25,0.85) {\small X 射线};
  \node at (12.75,0.85) {\small 伽马射线};
  \draw[->] (8.81,-0.15) -- (8.81,-0.6);
  \node[anchor=north] at (8.81,-0.6) {\small 可见光只有这一窄条};
  \foreach \x/\l in {0/$10^{3}$, 4.5/$10^{9}$, 6.75/$10^{12}$, 10.5/$10^{17}$, 12/$10^{19}$, 13.5/$10^{21}$}
    \node[below] at (\x,0) {\scriptsize \l};
  \node[below] at (6.75,-0.55) {\scriptsize 频率 / Hz};
\end{tikzpicture}
\end{center}

\begin{mainline}
认识一下家庭成员。无线电波频率最低、波长最长，从广播、对讲机到卫星电视都靠它，波长在米到千米的量级；微波短一些，Wi-Fi 工作在约 $2.4\times10^{9}$ 赫兹和 $5\times10^{9}$ 赫兹，微波炉用 $2.45\times10^{9}$ 赫兹，波长十几厘米；红外线更短，遥控器的发光管、夜视仪、你脸上向四周散发的热辐射都在这一段。顺带一个生活小魔术：电视遥控器的红外发光管人眼看不见，用手机摄像头对着它按下按键，屏幕上却能看见一亮一亮——摄像头的硅片对近红外还有响应，你的眼睛没有。可见光只占图上那条窄缝，波长约 380 到 760 纳米；紫外线能打断分子里的化学键，所以能晒黑皮肤、杀死细菌，消毒柜里那盏发蓝紫光的灯靠的就是它；X 射线穿过软组织、被骨头挡住，于是有了医院的片子；伽马射线来自原子核内部的变化，频率最高、穿透最强，医院放疗和考古测年都会用到它。

注意这张图的两端都没有“到头”：比伽马射线频率更高的电磁波原则上存在，只是越来越难产生和探测；比无线电波频率更低的也一样。这不是七种不同的东西排队站好，而是同一种东西的一把连续的尺子——人类只是按“它撞到我们时会发生什么”给各段起了名字。

这里插一项工程应用：雷达。雷达站发出一束微波脉冲，脉冲撞上飞机弹回来，测量来回的时间差乘以 $c$ 再除以二，就是距离；交警手里的测速仪用的也是同一招，只是还多看一眼回波频率的微小移动。导航卫星反过来用：几颗卫星不停地广播自己的位置和发信时刻，你的手机比较几路信号到达的时间差，就反推出自己在地球表面的位置——整套系统的根基，就是“电磁波以恒定速度 $c$ 走直线”这一条。巨型射电望远镜则什么都不发射，只竖一口大锅，静静接收宇宙深处天体发来的无线电波——人类对宇宙的大部分认识，都是顺着这家族的不同成员“听”来的。
\end{mainline}

\begin{mathbridge}
为什么波速恰好是 $1/\sqrt{\varepsilon_0\mu_0}$？用一个一维简化模型就能看到骨架。

设电场只有 $y$ 分量，只沿 $x$ 方向变化，记作 $E(x,t)$；磁场只有 $z$ 分量，记作 $B(x,t)$。取一个细窄的矩形回路平放在 $x$-$y$ 平面上，对它用第 23 章的法拉第定律：回路一条长边上的电场与另一条长边上的电场之差，等于穿过回路的磁通变化率。把回路收到无限窄，得到
\[
  \frac{\partial E}{\partial x}=-\frac{\partial B}{\partial t}\,.
\]
再把矩形回路平放到 $x$-$z$ 平面，对它用带位移电流补丁的安培定律（真空、无电流），同样取极限，得到
\[
  -\frac{\partial B}{\partial x}=\varepsilon_0\mu_0\,\frac{\partial E}{\partial t}\,.
\]
对第一式关于 $x$ 再求一次变化率，代入第二式，两式合流：
\[
  \frac{\partial^2 E}{\partial x^2}=\varepsilon_0\mu_0\,\frac{\partial^2 E}{\partial t^2}\,.
\]
拿它和第 17 章波动方程的标准样子对比——某量的空间二阶变化率等于时间二阶变化率除以波速平方——立刻读出
\[
  c=\frac{1}{\sqrt{\varepsilon_0\mu_0}}\,.
\]
同样的两步操作对 $B$ 再做一遍，会得到一模一样的方程和同一个速度。电场波和磁场波被方程锁成同一列队伍，谁也快不了，谁也慢不了。
\end{mathbridge}

\begin{challenge}
三维情形不需要新的物理，只需要向量微积分。在真空无源区域，对法拉第定律 $\nabla\times\bm E=-\partial\bm B/\partial t$ 两边再取一次旋度，用恒等式 $\nabla\times(\nabla\times\bm E)=\nabla(\nabla\cdot\bm E)-\nabla^2\bm E$，而真空里 $\nabla\cdot\bm E=0$，再代入安培—麦克斯韦定律，得到
\[
  \nabla^2\bm E=\varepsilon_0\mu_0\,\frac{\partial^2\bm E}{\partial t^2}\,,
\]
磁场满足完全相同的方程。这表示 $\bm E$ 的每一个直角坐标分量都是一个独立的三维波动方程。

横波性也从方程里直接掉出来：对平面波 $\bm E=\bm E_0\cos(\bm k\cdot\bm r-\omega t)$，条件 $\nabla\cdot\bm E=0$ 给出 $\bm k\cdot\bm E_0=0$，即振动方向垂直于传播方向；$\nabla\cdot\bm B=0$ 对磁场给出同样结论。再注意，这个方程里从头到尾没有出现任何参考系——$c$ 是“相对于谁”的速度？把这个问题记下来，它是第 38 章的入口。
\end{challenge}

\step{模型的边界}

\begin{mainline}
电磁波模型强得惊人，但它有自己的边界。

边界一：近处不是波。紧贴天线、变压器、充电线圈的地方，场的形态复杂，随距离迅速衰减，大部分能量绕一圈又回到电路里，并没有“飞走”。只有远离波源几个波长之后，场才整理成自由传播的电磁波。所以“一切变化的场都是电磁波”的说法过头了——只有被源真正“甩出去”的那部分才是。

边界二：介质里速度会变。在玻璃、水、空气里，电场要不断摇动材料里的电荷，波速降成 $c/n$，$n$ 是折射率。光在水里只有约 $0.75\,c$。真空中那个 $c$ 才是写进麦克斯韦方程的常数（介质中的故事见第 36 章）。

边界三：极弱的光会让这个模型破产。把光一再调暗，探测器会开始“一颗、一颗”地响应，而不是接收到一份连续变弱的波。平滑的电磁场图像是极好的模型，但不是最后的词——光还藏着颗粒性的一面，这是第 42 章开始的量子故事。这里要分清层次：“光在大量光子的情形下表现得像经典电磁波”是实验事实；“电磁波由麦克斯韦方程描述”是被反复验证的数学模型；“光就是场在振动”是至今最好用的物理解释，三者不是一回事。

边界四：那个没有回答的问题。$c=1/\sqrt{\varepsilon_0\mu_0}$ 是相对谁的速度？水波的速度相对水，声波的速度相对空气，而麦克斯韦方程里根本没有“以太相对谁静止”这一项。这个缺口不是小毛病——它会一路通向狭义相对论（第 38 章）。

典型误用也列三个。其一，把电磁波想象成“被空气吹散的涟漪”——恰恰相反，空气只会添乱，电磁波在真空里走得最干净。其二，流行说法“微波炉选在水分子的共振频率上”并不准确：$2.45\times10^{9}$ 赫兹并不是水的什么尖锐共振峰，而是加热效率、穿透深度和频段分配之间的工程折中。其三，一听见“辐射”两个字就紧张。本章的模型恰好给出了区分的关键：频率。可见光以下的电磁波，单个波包携带的能量不足以打断化学键，对身体的途径主要是“加热”；紫外线以上才有打断键、损伤分子的本领。“辐射”不是一个危险标签，而是一把频率尺子上的不同刻度。
\end{mainline}

\begin{mainline}
还有两个由模型直接兑现的重要结论。

一是天线的原理。发射天线就是一根让电荷来回加速的导线：匀速运动的电荷只是拖着自己的场一起走，加速的电荷才会把场“甩出去”成为自由波。天线尺寸和波长相当才甩得高效，所以中波广播台要竖起上百米的铁塔（波长约三百米），收音机拉杆天线几十厘米（对应波长约三米的调频波段），手机工作在约 $2\times10^{9}$ 赫兹、波长约十五厘米，几厘米的天线就能藏进机身。接收天线里没有“吸信号”的神秘力量：路过的电场推着金属里的自由电子来回跑，形成微弱电流，放大器把它读出来。发射与接收是同一条物理规律的两个方向。这也解释了手机上的“信号格”：格数本质上是天线收到的电场推动电子的强弱，基站远、墙体厚、手握住了天线的位置，推动就变弱。电梯里信号消失，则是因为金属轿厢成了一整套屏蔽壳——和你用铝箔包住收音机是同一回事。

\begin{center}
\begin{tikzpicture}[>=Stealth]
  \draw[very thick] (-2.2,0) -- (-0.25,0);
  \draw[very thick] (0.25,0) -- (2.2,0);
  \fill (0,0) circle (1.6pt);
  \draw[<->] (-1.3,0.4) -- (1.3,0.4) node[midway,above]{\small 电荷来回振荡};
  \draw[decorate,decoration={snake,amplitude=1.5pt,segment length=7pt},->] (2.5,0) -- (4.9,0) node[right]{\small 电磁波};
  \node[below] at (0,-0.25) {\small 发射天线};
\end{tikzpicture}
\end{center}

二是偏振与横波。电场的振动方向垂直于传播方向，而这个垂直方向上还可以有不同的取向：上下振、左右振、斜着振——这就是偏振。纵波做不到这一点：声波沿传播方向振动，绕传播轴转一圈哪个方向都一样，谈不上偏振。所以“光可以偏振”这件事本身，就是光是横波的硬证据。太阳光是各种取向混杂的非偏振光；偏振片内部像一排极细的栅栏，只放行某一个方向的电场分量。两片偏振片取向垂直时，第一个方向的分量全被第二片挡住，于是全黑——动手实验里的黑暗就是这么来的。这个比喻要交代边界：偏振片像栅栏的地方在于，它按方向筛选振动；它不像栅栏的地方在于，被“拦住”的那一半电场分量并没有排队从缝里挤过去，而是被材料吸收变成了热。钓鱼的人戴偏振太阳镜，利用的正是同一原理：水面反射的光以水平方向的振动为主，竖直放行的镜片把这层刺眼的反光大部分滤掉，水下的鱼就显出来了。那斜插的第三片为什么又能“救活”一点光？因为每一片不是简单地“减光”，而是把入射电场投影到自己的放行方向上：斜四十五度的中间片从第一片的输出里取出一个斜向分量，最后一片再从这个斜向分量里取出自己的分量，两次投影之后还剩下一部分。挡光板越挡越黑，投影器却可能“越挡越亮”——这就是直觉在这里出错的地方（定量的马吕斯定律见第 36 章）。
\end{mainline}

\step{回头解释开场现象}

\begin{mainline}
现在回到晒太阳。太阳的核心和表面，到处是被热运动不断加速的带电粒子，它们每时每刻都在把一团团电磁波“甩”进空间。其中一小撮朝地球飞来，用约八分二十秒跑完一亿五千万千米的真空。一路上没有任何东西在“运送”它——它只是按麦克斯韦方程的要求，自己维持自己，向前平移。到达你的皮肤时，电场推着分子里的电荷来回抖动，抖动积累的能量变成热（这笔账留给第 26 到 28 章）。你眼皮感到的那一点暖意，是八分钟前离开太阳的一份电场与磁场，此刻正在你脸上推电荷。

收音机、Wi-Fi、遥控器、X 光片，全是同一件事的不同频率版本。动手实验里的“咔哒”声也清楚了：电池火花甩出的宽带电磁波里，恰好有落在收音机能接收频段内的成分；而铝箔包住收音机后，外来电场推动铝箔里的自由电子重新分布，在箔内产生一个几乎抵消外来场的反向场——波不是被“吸掉”，而是被导体自己长出来的场顶了回去。

开场悬念的答案：穿越真空的不是粒子流，也不是神秘介质里的涟漪，而是场本身。

开灯“立刻”亮的悬念也一并解决。光从台灯到你眼睛只要几纳秒，以人的任何感官都无从分辨，所以“立刻”只是“快到超出日常经验”的同义词。但一旦距离拉到天文尺度，这个有限的速度就变成了看得见的事实：我们看见的月亮是一秒多以前的月亮，太阳是八分钟以前的太阳，仙女座星系是两百五十万年以前的仙女座星系。抬头看星空，其实是在看一部由光速亲手剪辑的“历史纪录片”——这也将是第 38 章重新审查时间与空间概念的出发点。
\end{mainline}

\step{脑洞实验与挑战题}

\begin{thought}[给太阳盖上盖子]
假如太阳此刻突然熄灭，我们会立刻知道吗？不会。最后一份阳光还在路上跑八分二十秒，这段时间里地球一切如常，没人能提前收到任何消息。更妙的是：连“太阳没了”的引力变化也快不了——引力的改变同样以光速传播（这是第 41 章的预告）。

把尺度推到极端：给 4.2 光年外的比邻星打一个电话，你说“喂”到听见对方回“喂”，要等八年多；银河系直径约十万光年，开一场银河视频会议，每句话都要延迟十万年。电磁波是宇宙里最快的信使，而宇宙大到连最快的信使也显得步履蹒跚。
\end{thought}

先做一个能用真实数字的估算，再来三道挑战题。

\textbf{估算：用一炉奶酪称出光速。}让微波炉的转盘停转（或架住不动），平铺一大盘芝士片或平放一条巧克力，低火加热十几秒，表面会出现一串融化了的“热点”，热点之间是凉的。热点对应腔内驻波的波腹，相邻热点间距是半个波长（驻波见第 17 章）。量一量，约为 6 厘米，于是波长 $\lambda\approx 12$ 厘米。炉门铭牌上标着频率 $f=2450$ 兆赫兹 $=2.45\times10^{9}$ 赫兹。于是
\[
  c=\lambda f\approx 0.12\times 2.45\times10^{9}\approx 2.9\times10^{8}\ \text{m/s}\,,
\]
与公认值只差约 $3\%$。用一炉奶酪测出一个宇宙基本常数——这就是“光是电磁波”最直接的日常证据。

\begin{puzzles}
\item 微波炉的玻璃门上贴着一片金属网，网孔只有一两毫米。你看得清里面的食物，微波却漏不出来。为什么？如果把网换成一条长 20 厘米、宽 1 毫米的细缝，还能放心吗？\\
思路提示：能不能穿过一个孔，取决于波长和孔的“最大尺寸”之比，而不取决于材料看起来透不透明。再想想细长缝的长边已经比微波波长长，会发生什么。

\item 把正在通话的手机放进盖紧的金属饭盒，信号断了；放进塑料盒，照常通话。可有人把手机放进（不通电的）旧微波炉里测试，电话有时仍能打通。解释这些现象。\\
思路提示：屏蔽的本质是导体上的感应电荷与电流产生反向场，把外来场抵消掉；只要存在与波长可比的缝隙或接触不良的接点，这条“抵消生产线”就会漏。手机波长约十几到几十厘米，旧炉门的封条未必处处密合。

\item 汽车开进长隧道，常常是调频 FM（波长约 3 米）还剩一点信号，中波 AM（波长约 300 米）却先消失。可平时又说“长波擅长绕过山和建筑”。矛盾吗？\\
思路提示：把隧道口看成一个“孔径”：孔径远大于波长时，波容易径直进出；孔径远小于波长时，波几乎进不去。绕射擅长的是“绕过障碍”，隧道考验的是“钻进洞口”，这是两回事。本题开放，重在用“波长与障碍尺寸的比较”组织推理。

\item 老式收音机收中波时，把拉杆天线转个方向，声音会明显变化；收调频时影响却小一些。猜一猜原因，再想想为什么“天线拉长一点信号就好”有时是错觉。\\
思路提示：天线里的电子是被电场的方向推着跑的，天线顺着电场方向摆，收到的推动最大；垂直摆就几乎收不到。中波台远，传播途中电场方向比较固定；调频信号在城市里被楼宇多次反射，到达时方向已经乱了。天线“变长”真正改变的，是它尺寸与波长的匹配程度。
\end{puzzles}
